\section{Conclusion}
\label{sec:conclusion}

Sparse-PMU attribution is a comparison of physical explanations over time. The implemented load/line bank predicts measurements from intervention-conditioned DAE histories, retains recovery states, and mixes fixed network hypotheses. $M$ records the resulting source masses; $\Delta M$ separates within-model evidence from changes in model weight. These operations make the decision traceable without turning standard Bayesian or filtering identities into novelty claims.

The completed development evidence shows both utility and a limit. Physical state recovery can be accurate while weak event sources remain ambiguous. In causal H04 development, the five-model reference obtains \ContMixTimely{}/\ContMixTotal{} timely joint decisions with \ContMixExtras{} extras; nominal improves timeliness at a higher alarm burden. The family and temporal breakdowns explain that result, including the BUS31 delay and the negative BUS16 admission experiment. Final performance claims depend on the frozen confirmatory results. No broader all-family, unseen-source or field claim follows from these tests.

\section*{Supplementary Material}
The supplement preserves detailed proofs, electrical and instrumental conventions, the original batch estimator and joint-estimator proposals, expanded component studies, and complete evidence contracts. The pre-revision manuscript is archived in the repository. Empirical data and inference baselines are unchanged.
